\section{Fundamentos te\'oricos}
\label{sec:fundamentos}

Esta secci\'on resume el marco te\'orico que sustenta el m\'etodo MASW, elaborado a partir
del estudio sistem\'atico del texto de referencia de Foti \emph{et al.}~\cite{Foti2014}
y de la literatura complementaria. El material fue organizado durante el proyecto como
un \emph{vault} de conocimiento en Obsidian, con los ocho cap\'itulos del libro resumidos
y m\'as de 140 conceptos at\'omicos interconectados (ver Secci\'on~\ref{sec:cronologia}).

\subsection{Comportamiento del suelo y elasticidad a peque\~nas deformaciones}

Los suelos son materiales particulados, multif\'asicos y discontinuos, cuyo comportamiento
mec\'anico es inherentemente no lineal e irreversible. Sin embargo, la elasticidad lineal
provee un marco consistente para interpretar los ensayos s\'ismicos, porque los niveles
de deformaci\'on que producen las ondas s\'ismicas son extremadamente peque\~nos: t\'ipicamente
$\gamma<10^{-5}$, muy por debajo del umbral no lineal del suelo ($\gamma\approx10^{-4}$
a $10^{-3}$). A estas deformaciones la curva esfuerzo--deformaci\'on es esencialmente
lineal y el m\'odulo de corte alcanza su valor m\'aximo $G_{\max}$.

La respuesta el\'astica lineal de un s\'olido is\'otropo queda completamente descrita por
dos constantes el\'asticas ---por ejemplo las constantes de Lam\'e $\lambda$ y
$\mu\equiv G$---, y las velocidades s\'ismicas se expresan como

\begin{equation}
V_P=\sqrt{\frac{\lambda+2G}{\rho}},\qquad V_S=\sqrt{\frac{G}{\rho}},
\label{eq:velocidades}
\end{equation}

donde $\rho$ es la densidad. Esta relaci\'on directa entre velocidad y m\'odulo el\'astico
es la base f\'isica de toda la s\'ismica aplicada.

Una observaci\'on relevante para la elecci\'on del m\'etodo: en suelos saturados la
compresibilidad del fluido intersticial domina la respuesta volum\'etrica, haciendo que
$V_P\to\SI{1500}{\meter\per\second}$ (la velocidad del sonido en agua) sin relaci\'on con
la rigidez del esqueleto s\'olido. $V_S$, en cambio, refleja \'unicamente la rigidez de
corte del esqueleto y es insensible al contenido de agua. Por eso $V_S$ ---y no
$V_P$--- es el par\'ametro objetivo de la caracterizaci\'on geot\'ecnica.

La disipaci\'on de energ\'ia observada incluso a deformaciones muy peque\~nas se modela
extendiendo el marco a viscoelasticidad lineal, reemplazando el m\'odulo real por el
m\'odulo complejo $G^{*}=G(1+2iD_S)$, con $D_S$ la raz\'on de amortiguamiento material.
Esto introduce un n\'umero de onda complejo $k^{*}=k+i\alpha$, cuya parte real controla
la velocidad de fase y cuya parte imaginaria controla el decaimiento espacial de
amplitud.

\subsection{Ondas de cuerpo y ondas superficiales}

En un medio el\'astico is\'otropo infinito se propagan dos tipos de ondas de cuerpo: las
ondas P (compresionales) y las ondas S (de corte). La presencia de una superficie libre
genera, adem\'as, ondas superficiales: las ondas Rayleigh (movimiento el\'iptico retr\'ogrado
en el plano vertical) y las ondas Love (movimiento horizontal transversal, que requieren
estratificaci\'on).

Para los m\'etodos de caracterizaci\'on desde superficie, las ondas Rayleigh son las
relevantes por dos motivos cuantitativos. Primero, la partici\'on de energ\'ia: a distancias
de una o dos longitudes de onda de una fuente puntual vertical, las ondas superficiales
concentran aproximadamente el \SI{67}{\percent} de la energ\'ia total radiada. Segundo,
el decaimiento geom\'etrico: las ondas superficiales decaen como $r^{-0.5}$ frente a
$r^{-2}$ de las ondas de cuerpo en superficie. La combinaci\'on de ambos efectos hace
que, a distancia suficiente de la fuente, el registro de un ge\'ofono vertical est\'e
dominado por el tren Rayleigh ---el llamado \emph{ground roll}, que en s\'ismica de
reflexi\'on se considera ruido y que aqu\'i es precisamente la se\~nal de inter\'es.

\subsection{Dispersi\'on geom\'etrica: el principio del m\'etodo}

En un semiespacio homog\'eneo la velocidad de fase Rayleigh $V_R$ es independiente de la
frecuencia y vale aproximadamente $0{,}9\,V_S$. La dispersi\'on aparece cuando el medio
es verticalmente heterog\'eneo.

El mecanismo es geom\'etrico: la profundidad de penetraci\'on de una onda Rayleigh es
proporcional a su longitud de onda $\lambda$. Componentes de baja frecuencia (longitud
de onda larga) involucran un espesor grande de suelo y por lo tanto su velocidad de
fase queda determinada por un promedio ponderado que incluye las capas profundas;
componentes de alta frecuencia (longitud de onda corta) muestrean s\'olo las capas
superficiales. En consecuencia, en un perfil normalmente dispersivo (rigidez creciente
con la profundidad) la velocidad de fase crece al disminuir la frecuencia.

La funci\'on $V_R(f)$ ---la \textbf{curva de dispersi\'on}--- es entonces una firma del
perfil $V_S(z)$. El m\'etodo MASW consiste en medir experimentalmente esa curva y
resolver el problema inverso para recuperar el perfil.

\subsection{Modos de propagaci\'on}

En un medio estratificado la relaci\'on de dispersi\'on admite m\'ultiples soluciones para
cada frecuencia: el modo fundamental y los modos superiores. Cada modo tiene su propia
curva $V_R(f)$ y su propia distribuci\'on de energ\'ia con la profundidad.

Esto tiene dos consecuencias pr\'acticas importantes. La primera es que la curva
experimental medida no siempre corresponde al modo fundamental: en perfiles con
inversiones de velocidad (una capa r\'igida sobre una blanda) o con contrastes fuertes,
los modos superiores pueden dominar la energ\'ia en ciertos rangos de frecuencia, y lo
que se mide es en realidad una \emph{curva aparente} o \emph{efectiva} que salta entre
modos. Interpretarla como fundamental produce un perfil $V_S$ err\'oneo. La segunda es
que, cuando los modos superiores pueden identificarse por separado, aportan informaci\'on
adicional que reduce la no-unicidad de la inversi\'on. Por este motivo el
\emph{pipeline} desarrollado en este proyecto implementa extracci\'on e inversi\'on
multimodal (Secci\'on~\ref{sec:arq-pipeline}).

\subsection{El flujo de trabajo: adquisici\'on, procesamiento e inversi\'on}

La metodolog\'ia est\'andar de los m\'etodos de ondas superficiales sigue tres pasos
encadenados:

\begin{enumerate}[itemsep=0.15em]
    \item \textbf{Adquisici\'on}: registro del movimiento vertical de la superficie
    mediante un arreglo lineal de ge\'ofonos, excitado por una fuente impulsiva
    (m\'etodo activo) o aprovechando ruido ambiental (m\'etodo pasivo).
    \item \textbf{Procesamiento}: transformaci\'on del registro del dominio
    espacio--tiempo $(x,t)$ a un dominio donde la dispersi\'on sea observable
    ---frecuencia--n\'umero de onda $(f,k)$, frecuencia--velocidad de fase $(f,V)$ o
    frecuencia--lentitud--- y extracci\'on de la curva de dispersi\'on como el lugar
    geom\'etrico de los m\'aximos espectrales (\emph{picking}).
    \item \textbf{Inversi\'on}: resoluci\'on del problema inverso no lineal para obtener
    el perfil $V_S(z)$ que mejor reproduce la curva experimental.
\end{enumerate}

\subsubsection{Limitaciones de la adquisici\'on}

El dise\~no del arreglo impone l\'imites duros al rango de frecuencias recuperable, y estos
l\'imites condicionan directamente las especificaciones del sistema electr\'onico:

\begin{itemize}[itemsep=0.15em]
    \item \textbf{L\'imite superior de frecuencia}: determinado por el espaciamiento
    entre receptores $\Delta x$ v\'ia \emph{aliasing} espacial ($\lambda_{\min}=2\Delta x$),
    y por la atenuaci\'on de las componentes de alta frecuencia en los ge\'ofonos lejanos.
    \item \textbf{L\'imite inferior de frecuencia}: determinado por la longitud total
    del arreglo $L$, que fija la longitud de onda m\'axima resoluble y, por lo tanto, la
    profundidad m\'axima de investigaci\'on (regla pr\'actica: $z_{\max}\approx L/2$ a $L/3$).
    \item \textbf{Efecto de campo cercano}: a distancias de la fuente menores a
    aproximadamente media longitud de onda, el campo de ondas no est\'a a\'un dominado por
    la onda superficial plana y la velocidad de fase medida est\'a sesgada. Esto impone
    una distancia m\'inima fuente--primer ge\'ofono.
\end{itemize}

Un objetivo de \SI{50}{\meter} de profundidad implica entonces arreglos de longitud
considerable y, sobre todo, capacidad de recuperar componentes de baja frecuencia con
buena relaci\'on se\~nal--ruido. Este es exactamente el requisito que tensiona la
instrumentaci\'on, y el que motiva tanto el trabajo de compensaci\'on de la respuesta del
ge\'ofono como el de calibraci\'on del \emph{front-end} descritos m\'as adelante.

\subsubsection{El problema inverso}

La inversi\'on es un problema mal condicionado y no \'unico: distintos perfiles $V_S(z)$
pueden producir curvas de dispersi\'on pr\'acticamente indistinguibles dentro del error
experimental. Las estrategias habituales combinan parametrizaci\'on en capas,
regularizaci\'on, informaci\'on a priori y algoritmos de b\'usqueda global. La inversi\'on
requiere adem\'as fijar par\'ametros auxiliares ---densidad y raz\'on de Poisson por capa---
que tienen influencia limitada pero no nula sobre la curva; en particular, ignorar el
nivel fre\'atico equivale a asumir una raz\'on de Poisson incorrecta para las capas
saturadas (donde el comportamiento no drenado lleva $\nu\to0{,}5$), con consecuencias
significativas sobre el perfil resultante.

\subsection{Del m\'etodo a la instrumentaci\'on}
\label{sec:fund-geofono}

Los l\'imites de adquisici\'on enunciados arriba se traducen en requisitos concretos sobre
la electr\'onica: recuperar contenido de baja frecuencia con buena relaci\'on
se\~nal--ruido, y garantizar que todos los canales del arreglo compartan exactamente la
misma referencia temporal. El transductor que impone esos requisitos ---el ge\'ofono--- se
trata en detalle en la Secci\'on~\ref{sec:geofono}, y las topolog\'ias de acondicionamiento
estudiadas para satisfacerlos, en la Secci\'on~\ref{sec:topologias}.

\nota{Para la versi\'on de 15 p\'aginas esta secci\'on deber\'ia bajar a ~2 p\'aginas:
conservar \S2.1 (reducido a la Ec.~\ref{eq:velocidades} y el argumento de $V_S$ vs
$V_P$), \S2.3 completo (es el principio del m\'etodo) y \S2.5.1 (porque justifica las
especificaciones del hardware). \S2.2 y \S2.4 pueden comprimirse a un p\'arrafo cada uno.}
